% 与主文件同步的模板第2、3部分；使用主文件的宏定义与TikZ设置。
\subsection{项目的研究内容、研究目标及拟解决的关键科学问题}

\subsubsection{研究内容}

本项目面向柔性桁架上机器人的精确运动与模块组装，以机器人动力学、几何阻抗、运动优化与接触几何为基础，采用理论分析、数值仿真与地面试验相结合的研究手段，研究桁架振动条件下的耦合动力学与运动控制、目标位姿不确定条件下的接触对准与柔顺组装方法，并开展地面试验验证。具体研究内容如下：

\topic{（1）柔性桁架上机械臂的耦合动力学与运动控制}

机器人依附柔性桁架作业时，机械臂运动通过安装端反作用力与力矩激发桁架振动，桁架振动又通过机器人基座的位置与姿态变化影响末端位姿与动态响应。针对上述双向耦合对末端运动与接触响应的影响，从安装界面的运动与受力出发，建立面向实时控制的倒置动力学描述，研究几何阻抗与动力学优化相结合的运动控制方法，为振动环境下的精确运动与接触作业提供基础。主要研究内容如下：
\begin{enumerate}[label=\arabic*)]
 \item 研究基座、关节与末端之间的运动耦合及力和力矩传递关系，通过安装端运动与受力表征柔性桁架与机械臂之间的耦合作用，建立以作业端为浮动基的倒置动力学模型，分析基座振动、末端负载及机械臂运动状态对末端响应的影响规律。
 \item 面向位姿保持与轨迹跟踪任务，研究几何阻抗与动力学优化控制方法，统一描述末端位置与姿态响应，在倒置动力学框架下协调关节运动与力矩分配，并计入关节运动范围、驱动力矩及其变化率限制，实现基座振动条件下的精确运动控制。
 \item 研究安装端作用力与力矩的建模偏差及驱动约束对实际阻抗响应的影响，揭示末端运动精度、接触柔顺性与安装端反作用之间的关联，明确不同基座振动、运动速度及末端负载条件下的控制参数选取依据与方法适用范围。
\end{enumerate}

\topic{（2）目标位姿不确定条件下的接触对准与柔顺组装}

针对目标位姿估计偏差与基座运动共同作用下接口对准困难、接触力易增大的问题，研究接触几何对目标位姿偏差的约束作用，利用接触位姿与接触力修正组装目标及运动参考，并将修正后的参考输入运动控制过程，实现接口对准、插接与到位保持过程中的精度与柔顺性协调。主要研究内容如下：
\begin{enumerate}[label=\arabic*)]
 \item 研究典型接口的局部接触流形表征方法，分析接口相对位姿、接触状态与接触力之间的关联，明确不同接触信息对横向位置及姿态偏差的约束能力，为目标位姿修正提供几何依据。
 \item 考虑基座运动对接触位姿测量与几何判断的影响，研究统一参考系下的接触场配准与目标位姿估计方法，结合目标先验和接触反馈构造有界参考修正，实现横向及姿态纠偏与轴向接近的协同。
 \item 研究接触参考修正与实际阻抗响应的协调机制，分析目标估计偏差、运动执行误差及局部接触约束对接触力和安装端受力的影响，明确基座振动条件下的稳定组装条件与可完成对接的初始位姿偏差范围。
\end{enumerate}

\topic{（3）柔性桁架上机器人运动与组装的地面试验验证}

依托基座振动模拟与机器人作业实验平台，开展地面试验。平台由基座运动模拟装置、作业机械臂及同步测量系统组成。其中，基座运动模拟装置采用机械臂实现，其末端与作业机械臂基座固定连接，通过施加可重复的平移与转动激励，模拟柔性桁架振动引起的安装端运动；作业机械臂执行运动与组装任务。围绕不同振动条件、末端负载及初始位姿偏差，开展运动控制与典型接口组装的分项及集成试验，验证所提动力学描述与控制方法，明确其适用范围。主要研究内容如下：
\begin{enumerate}[label=\arabic*)]
 \item 完善悬挂连接、典型接口与同步测量系统，开展基座与末端位姿测量系统、安装端与接触端力和力矩传感器的标定，设置不同幅值、频率和方向的基座运动激励，建立可重复的基座振动模拟与试验评价方法。
 \item 开展位姿保持与轨迹跟踪试验，通过不同基座振动、运动速度及末端负载条件下的数值仿真与实测响应对照，校核安装端作用力与力矩的动力学描述，评价所提运动控制方法的运动精度、力矩约束满足情况与实时执行能力。
 \item 开展不同初始位姿偏差下的典型接口对准、插接及到位保持试验，对比固定参考、接触场配准及接触参考修正的作用，综合评价组装到位误差、接触力与安装端反作用，完成振动环境下机器人运动与组装的地面原理验证。
\end{enumerate}

\subsubsection{研究目标}

面向大型空间结构在轨组装中柔性桁架上机器人的精确运动与柔顺操作需求，建立基于安装界面运动与受力的耦合动力学描述，揭示桁架振动、关节驱动约束与末端运动响应之间的关系，形成几何阻抗与动力学优化相结合的运动控制方法；阐明接触几何对目标位姿偏差的约束作用，建立接触场配准与有界参考修正相结合的柔顺组装方法，实现基座振动及目标位姿不确定条件下的接口对准、插接与到位保持；依托基座振动模拟与机器人作业实验平台，验证所提动力学描述与控制方法的有效性，明确其在不同振动、末端负载及初始位姿偏差条件下的适用范围，形成理论分析、数值仿真与地面试验相互印证的研究框架。

\subsubsection{拟解决的关键科学问题}

\topic{（1）机器人与柔性桁架耦合作用下受约束运动与阻抗响应的形成机制}

柔性桁架振动通过安装端运动与受力改变机械臂的动态响应，机械臂驱动又通过安装端反作用影响柔性桁架的振动响应，使作业端运动与安装端受力相互关联。面向实时控制，需要保留影响末端响应的主要耦合信息，而安装端作用力与力矩的近似描述会带来动力学偏差；关节运动范围与驱动力矩限制还可能使期望运动无法完全实现，导致实际阻抗响应偏离设计值。因此，安装端耦合作用如何经受约束力矩分配影响末端运动与柔顺响应，是本项目需要解决的关键科学问题。

需要揭示安装端运动与受力向末端位姿响应传递的规律，阐明安装端作用力与力矩的建模偏差、驱动约束及参考调整对实际阻抗特性的共同影响，建立末端运动精度、接触柔顺性与安装端反作用之间的联系，为柔性桁架振动条件下的受约束运动控制及参数选取提供依据。

\topic{（2）柔性桁架振动条件下接触几何修正与阻抗响应的协调机制}

对接时的接触位姿和接触力同时受目标偏差、基座运动与机器人柔顺响应影响，较小的接触力不一定对应正确的装配位置，局部接触数据也未必能充分约束所有位姿方向。利用接触观测更新目标或运动参考，又会改变接触状态、关节作用及安装端受力，使几何估计与运动控制相互影响。

需要明确接触信息约束目标偏差的能力及可信程度，揭示基座振动条件下参考变化经实际阻抗控制传递至接触力和安装端受力的规律，使参考调整既能利用接触反馈纠偏，又保持对目标位姿的约束，从而协调振动扰动下的对准精度与柔顺性，实现典型工况下的稳定对接。

\subsection{拟采取的研究方案及可行性分析}

\subsubsection{总体研究思路与技术路线}

对应三项研究内容，按照“耦合动力学与运动控制、接触引导柔顺组装、地面试验验证”的顺序实施。首先以作业端为浮动基建立倒置机械臂动力学模型，通过实际安装端的作用力与力矩计入柔性桁架对机械臂的作用。在此框架下，结合几何阻抗设计与动力学优化，调节末端运动响应，协调关节运动与力矩分配；进一步利用接触流形的几何约束，通过接触场配准估计组装目标，并以有界参考修正实现柔顺对准与插接；最后利用地面实验平台模拟基座振动，开展分项对照与集成验证。技术路线如图\ref{fig:route}所示，各项方案分别说明如下。

\begin{figure}[htbp]
\centering
\begin{tikzpicture}[
 block/.style={draw=blue!45!black,fill=blue!3,rounded corners=3pt,
   text width=12.4cm,align=center,inner sep=8pt,font=\small},
 flow/.style={-{Stealth[length=2.5mm]},thick,draw=blue!50!black},
 node distance=0.55cm]
 \node[block,fill=blue!8] (need) {\textbf{柔性桁架上的精确运动与柔顺组装}\\[3pt]
 基座振动\qquad 关节驱动约束\qquad 目标位姿偏差};
 \node[block,below=of need] (control) {\textbf{研究内容一\quad 耦合动力学与运动控制}\\[3pt]
 倒置动力学\quad 几何阻抗\quad 动力学优化};
 \node[block,below=of control,draw=green!40!black,fill=green!3] (contact) {\textbf{研究内容二\quad 接触对准与柔顺组装}\\[3pt]
 接触流形表征\quad 接触场配准\quad 有界参考修正};
 \node[block,below=of contact,draw=violet!60!black,fill=violet!3] (test) {\textbf{研究内容三\quad 地面试验验证}\\[3pt]
 基座振动模拟与同步测量\\[2pt]
 位姿保持\quad 轨迹跟踪\quad 典型接口组装};
 \draw[flow] (need) -- (control);
 \draw[flow] (control) -- (contact);
 \draw[flow] (contact) -- (test);
\end{tikzpicture}
\caption{项目技术路线与验证关系}
\label{fig:route}
\end{figure}
\FloatBarrier

\subsubsection{倒置动力学建模与运动控制}

沿用以作业端为浮动基的倒置动力学模型，将实际安装基座作为倒置运动链的虚拟末端，通过该处的作用力与力矩描述柔性桁架对机械臂的作用。结合关节状态与作业端的绝对位姿、速度测量，确定倒置模型的运动状态，使基座运动对末端的影响能够通过实际运动反馈进入控制过程。在该建模框架下，结合几何阻抗设计与动力学优化：由几何阻抗关系给出作业端参考加速度，利用动力学二次规划（QP）协调末端运动、关节运动与力矩分配。安装端作用力与力矩的计算以现有静力平衡近似为基线，结合安装端运动、作用力与力矩数据校核其在动态工况下的适用性，分析安装端作用力与力矩的建模偏差对实际控制响应的影响。

在地面实验平台中，作业机械臂基座与基座运动模拟装置末端固定连接，该连接处对应模型中的实际安装端；作业机械臂末端对应倒置模型的浮动基。模拟装置提供安装端运动激励，作业机械臂的关节力矩由所提控制方法生成。结合实测安装端运动、作用力与力矩，校核安装端受力模型。由此建立在轨桁架与机器人的相互作用、地面基座激励与控制模型之间的对应关系。

以$T_b,T_d\in\SE$分别表示作业端的实际与期望位姿，$\bm V_b$和$\bm V_d$分别为对应的体速度。采用从实际位姿指向期望位姿的误差定义：
\begin{equation}
\begin{aligned}
 \widetilde T&=T_b^{-1}T_d,\qquad
 \bm\lambda=\operatorname{Log}(\widetilde T)^\vee,\\
 \widetilde{\bm V}
 &=\operatorname{Ad}_{\widetilde T}\bm V_d-\bm V_b
 =\operatorname{dexp}_{\bm\lambda}\dot{\bm\lambda}.
\end{aligned}
\label{eq:pose_error}
\end{equation}
其中，$\bm\lambda\in\mathbb R^6$为局部指数坐标，$\operatorname{Ad}$为群伴随映射，$\operatorname{dexp}_{\bm\lambda}$为与上述相对速度约定一致的指数映射微分。在李群上设计几何阻抗关系，并将其转换为指数坐标中的误差动态\cite{kim2025}：
\begin{equation}
\begin{aligned}
 \bm A\dot{\widetilde{\bm V}}+\bm D\widetilde{\bm V}
 +\operatorname{dexp}_{\bm\lambda}^{-\mathsf T}\bm K\bm\lambda
 &=\widetilde{\bm F},\\
 \bm A_\lambda\ddot{\bm\lambda}
 +\bm D_\lambda\dot{\bm\lambda}+\bm K_\lambda\bm\lambda
 &=\bm\gamma.
\end{aligned}
\label{eq:impedance}
\end{equation}
这里，$\bm A$、$\bm D$和$\bm K$分别为期望惯量、阻尼和刚度矩阵，$\widetilde{\bm F}$为按相对运动方向定义的外力旋量。记$\bm J_\lambda=\operatorname{dexp}_{\bm\lambda}$，则$\bm A_\lambda=\bm J_\lambda^{\mathsf T}\bm A\bm J_\lambda$、$\bm D_\lambda=\bm J_\lambda^{\mathsf T}\bm D\bm J_\lambda+\bm J_\lambda^{\mathsf T}\bm A\dot{\bm J}_\lambda$、$\bm K_\lambda=\bm K$、$\bm\gamma=\bm J_\lambda^{\mathsf T}\widetilde{\bm F}$。力与速度采用共轭映射，平移和转动按相应物理尺度设置参数。

由指数坐标中的阻抗关系计算参考加速度，再映射至倒置模型的浮动基：
\begin{equation}
\begin{aligned}
 \ddot{\bm\lambda}_{\rm ref}
 &=-\bm A_\lambda^{-1}\bm D_\lambda\dot{\bm\lambda}
   -\bm A_\lambda^{-1}\bm K_\lambda\bm\lambda
   +\bm A_\lambda^{-1}\bm\gamma,\\
 \dot{\bm V}_{\rm ref}
 &=\operatorname{Ad}_{\widetilde T}\dot{\bm V}_d
   -\operatorname{dexp}_{\bm\lambda}\ddot{\bm\lambda}_{\rm ref}
   +\operatorname{ad}_{\widetilde{\bm V}}\bm V_b
   -\dot{\bm J}_\lambda\dot{\bm\lambda}.
\end{aligned}
\label{eq:impedance_reference}
\end{equation}
其中，$\operatorname{ad}$为李代数伴随算子。浮动基速度与加速度分别记作$\dot{\bm q}_b$和$\ddot{\bm q}_b$，采用与$\bm V_b$及其导数一致的体坐标表示；因此，浮动基加速度参考为$\ddot{\bm q}_{b,\rm ref}=\dot{\bm V}_{\rm ref}$。由此将轨迹前馈与几何阻抗修正统一为动力学优化的参考输入。

按浮动基$b$与关节$j$划分倒置模型动力学，写为
\begin{equation}
 \begin{bmatrix}\bm M_{bb}&\bm M_{bj}\\\bm M_{jb}&\bm M_{jj}\end{bmatrix}
 \begin{bmatrix}\ddot{\bm q}_b\\\ddot{\bm q}_j\end{bmatrix}
 +\begin{bmatrix}\bm h_b\\\bm h_j\end{bmatrix}
 =\begin{bmatrix}\bm F_b\\\bm\tau_j\end{bmatrix},
 \qquad
 \begin{bmatrix}\bm h_b\\\bm h_j\end{bmatrix}
 =\bm C\begin{bmatrix}\dot{\bm q}_b\\\dot{\bm q}_j\end{bmatrix}
  +\bm G-\bm J_g^{\mathsf T}\bm F_e .
\label{eq:partitioned_dynamics}
\end{equation}
其中，$\bm M$、$\bm C$和$\bm G$分别为名义惯性矩阵、科氏及离心项矩阵和重力项；$\bm F_b$为作用在浮动基即实际作业端的接触力旋量，$\bm F_e$为作用在虚拟末端即实际安装端的作用力旋量，$\bm J_g$为其几何雅可比矩阵。自由运动时$\bm F_b=0$，接触阶段由测量或经校核的估计获得；$\bm F_e$以满足浮动基重力平衡的安装端静态作用力旋量为名义值，基座运动引起的安装端动态作用力旋量偏差结合实测数据进行评价。$\bm\tau_j$为待求解的命令等效关节力矩。

将动力学优化写成二次规划问题。令优化变量$\bm x=[\ddot{\bm q}_j^{\mathsf T},\bm\tau_j^{\mathsf T}]^{\mathsf T}$，以关节虚拟阻尼和驱动力矩正则为目标，将分块动力学作为线性等式约束：
\begin{equation}
\begin{aligned}
 \min_{\ddot{\bm q}_j,\,\bm\tau_j}\quad &
 w_{\ddot q}\|\ddot{\bm q}_j-\ddot{\bm q}_{j,\rm ref}\|^2
 +w_\tau\|\bm\tau_j-\bm\tau_{j,\rm ref}\|^2,\\
 \mathrm{s.t.}\quad &
 \begin{bmatrix}\bm M_{bj}&\bm 0\\\bm M_{jj}&-\bm I\end{bmatrix}
 \begin{bmatrix}\ddot{\bm q}_j\\\bm\tau_j\end{bmatrix}
 =\begin{bmatrix}
 \bm F_b-\bm h_b-\bm M_{bb}\ddot{\bm q}_{b,\rm ref}\\
 -\bm h_j-\bm M_{jb}\ddot{\bm q}_{b,\rm ref}
 \end{bmatrix},\\
 &-k_1\dot{\bm q}_j-k_0(\bm q_j-\bm q_{j,\min})
 \le\ddot{\bm q}_j,\\
 &\ddot{\bm q}_j\le-k_1\dot{\bm q}_j+k_0(\bm q_{j,\max}-\bm q_j),\\
 &\bm\tau_{j,\min}\le\bm\tau_j\le\bm\tau_{j,\max},\qquad
 |\bm\tau_j-\bm\tau_{j,k-1}|\le\Delta\bm\tau_{j,\max}.
\end{aligned}
\label{eq:qp}
\end{equation}
其中，$w_{\ddot q}$和$w_\tau$为正权重，$k_0,k_1>0$为关节边界约束参数，不等式按分量施加。关节参考取$\ddot{\bm q}_{j,\rm ref}=-\kappa\dot{\bm q}_j/\Delta t$，$\kappa>0$为虚拟阻尼系数，$\Delta t$为控制周期；取$\bm\tau_{j,\rm ref}=\bm 0$以减小关节驱动力矩。该式给出参考加速度可实现时的基本形式；当驱动或关节边界使其不可实现时，对浮动基加速度参考引入松弛并在目标中惩罚，保持动力学等式及驱动边界的一致性，分析参考放松对实际阻抗响应的影响。关节力矩及其变化量统一纳入约束；实际执行指令按机器人接口的重力约定换算。

在二次规划问题可行及状态估计误差有界的条件下，分析安装端作用力与力矩的建模偏差对位姿与速度误差动态的影响，研究局部闭环稳定与误差有界的条件。进一步考察约束激活时参考加速度松弛对实际阻抗响应的影响，结合基座运动、安装端受力及执行时序的实测数据校核扰动范围，并通过参数变化与约束激活试验检验分析结论的适用性。

\subsubsection{基于接触流形表征与参考修正的柔顺组装}

将典型接口在局部接触模式下允许的接触位姿集合视为接触流形，研究其对横向偏差和姿态偏差的几何约束。将接触位姿映射到统一的局部坐标，利用离线接触数据及物理方向信息构造接触场，以场的零值集合近似表征局部接触流形。方向信息作为场形状的软约束，以减轻接触采样不足的影响；在线阶段利用接触位姿相对于该场的残差估计目标：
\begin{equation}
 \widehat T_\star=\underset{T\in\mathcal U(T_v)}{\arg\min}
 \sum_i\rho\!\left(F(T^{-1}T_i)\right).
 \label{eq:registration}
\end{equation}
其中，$T_i$为统一参考系中的接触位姿，$T_v$为目标先验，$\mathcal U(T_v)$为先验附近允许修正的局部位姿集合，$F$表示离线接触场定义的带符号残差，$\rho$为鲁棒损失。该式提炼目标配准关系，具体采用横向和姿态修正，轴向接近单独调节。接触位姿及目标参考按同一时刻的基座运动统一坐标，避免将坐标变化误作目标偏差。

以配准目标生成名义对接参考，以$\bm c$表示局部横向和姿态修正量。将接触作用投影至这些修正方向，通过约束优化兼顾接触顺应与对配准目标的保持：
\begin{equation}
 \bm c_{k+1}=\underset{\bm c\in\mathcal C_k}{\arg\min}
 \left\{\frac{\|\bm c-\bm c_k\|_{\bm D_c}^2}{2\Delta t}
 +\frac12\bm c^{\mathsf T}\bm K_c\bm c
 -\bm f_{c,k}^{\mathsf T}(\bm c-\bm c_k)\right\}.
 \label{eq:capture}
\end{equation}
其中，$\bm f_{c,k}$为与修正坐标共轭的接触作用，$\bm D_c$为正定阻尼权重，$\bm K_c$依据配准敏感程度设置并正则化为正定的参考保持权重，$\mathcal C_k$限定累计偏移和单步变化。通过位姿映射将修正量转换为新的参考位姿$T_d$，并将其作为运动控制的输入。两处优化分别负责接触参考调整与关节力矩分配，保持功能和物理量定义清楚。

分析局部接触参考更新的储能变化，进一步考察下层跟踪误差与基座运动对接触力和对准的影响。参考子系统结论的整机适用性，以及有限转动和接触变化的影响，由包含基座运动的仿真与实验评价。

\subsubsection{地面试验系统与分项、集成验证}

依托现有地面实验平台及测量系统，通过基座运动模拟装置驱动作业机械臂基座产生可重复的平移或转动，模拟柔性桁架振动引起的安装端运动。作业机械臂在此条件下执行位姿保持、轨迹跟踪及典型接口组装任务。依据平台工作范围和预实验结果，按激励幅值、频率与方向设置多组代表性基座振动工况，通过更换末端配重及调整安装位置形成不同质量与惯量条件；以模拟装置保持静止时的同一安装状态作为补充基线，比较基座运动对作业性能的影响。

同步采集作业机械臂关节状态、实际基座运动、作业端位姿以及安装端、接触端的力和力矩，标定悬挂连接的空间外参及力和力矩传感器，并校准模拟装置与作业机械臂的执行时序。分别记录模拟装置的运动指令和作业机械臂基座的实测运动，以实测结果校核激励幅值、频率和相位；采用独立于控制反馈的测量结果评价末端精度与组装到位误差。安装端的力和力矩在统一坐标系下表达，并结合重力和静态受力标定分析其动态变化。

对应第一项研究内容，先通过位姿保持与圆轨迹跟踪试验校核倒置动力学模型与运动控制方法，对比数值仿真与实测响应，考察不同基座振动、末端负载及约束激活时的位姿误差、力矩边界与柔顺响应。记录位置均方根误差、旋转测地误差、关节力矩幅值与变化率，以及安装端力和力矩；同时记录控制计算耗时的分布、最大值及超出控制周期的比例，验证所提方法的实时执行能力。同一工况下，保持模拟装置的激励指令、作业机械臂的任务参考、执行时间与驱动约束一致，并核查实际基座运动的一致性，至少进行5次重复试验；参数对照在代表性工况开展。

对应第二项研究内容，在典型接口及不同初始偏差下，对比固定参考、配准后固定参考及配准结合接触修正，评价接触流形表征、目标配准和参考修正对组装对准的作用。由基座静止逐步过渡至模拟装置施加基座激励的工况，记录目标估计误差、接触力峰值、真实到位误差、作业完成情况及安装端力和力矩。目标安装方式在正式试验前固定，独立测量其在统一参考系中的位姿，按同时刻的末端与目标位姿评价相对对准误差；初始偏差与装配容差依据接口及预实验确定，调参和模型校核数据不计入独立验证。

对应第三项研究内容，完成上述分项试验及运动与组装的集成验证，围绕基座振动、末端负载与目标初始位姿误差设置对照，分析其对末端运动精度和接触对准的影响，明确各控制环节的贡献及方法适用范围。对于按给定运动施加的基座激励，以同等激励下的末端误差和接触响应评价控制效果，同时报告绝对误差与重复试验离散程度；安装端反作用力与力矩用于分析控制对柔性桁架的作用，其变化不直接等同于柔性结构振动幅值的变化。

\subsubsection{可行性分析}

\topic{（1）控制方法与实验接口具备直接基础}

现有位姿保持与圆轨迹跟踪实验已采用倒置动力学模型，并实现了几何阻抗与动力学优化控制。现有基座振动模拟与机器人作业实验平台可通过安装端运动施加振动扰动，支持几何微分、坐标映射与基座激励的校核，并可在此基础上开展安装端作用力与力矩的建模偏差及驱动约束影响分析，具备分步实施基础。

\topic{（2）接触研究与地面条件能够支撑方法衔接}

接触研究已形成配准与锚定捕获方法，具有仿真和真机配准回放基础。新方法真机捕获及其与力矩控制的融合拟由基座静止逐步过渡至模拟装置施加基座激励的工况。团队已具备地面实验平台及模块对接条件，需完善悬挂连接、接口夹具、安装端与接触测点，以及同步采集和系统集成。

\topic{（3）研究范围与实施难点相匹配}

试验由基座运动模拟装置施加运动激励，作业机械臂执行运动与组装任务，验证范围限定于预定振动条件和典型接口。利用实测基座运动与安装端受力校核动力学描述，结合末端误差和关节力矩调整阻抗参数与约束设置，利用先验与修正边界应对接触信息不足。各环节分别验证后再集成，使建模、控制和试验工作与两年开放课题条件相适应。

